
\documentclass[11pt]{article}

\usepackage[T1]{fontenc}
\usepackage{lmodern}
\usepackage[margin=1in]{geometry}
\usepackage{amsmath,amssymb,amsthm,mathtools}
\usepackage{booktabs,longtable,array}
\usepackage{enumitem}
\usepackage{microtype}
\usepackage{xcolor}
\usepackage{hyperref}
\usepackage{xurl}

\hypersetup{colorlinks=true,linkcolor=blue,citecolor=blue,urlcolor=blue}
\urlstyle{same}

\newtheorem{theorem}{Theorem}[section]
\newtheorem{corollary}[theorem]{Corollary}
\newtheorem{proposition}[theorem]{Proposition}
\newtheorem{lemma}[theorem]{Lemma}
\newtheorem{remark}[theorem]{Remark}
\newtheorem{definition}[theorem]{Definition}
\newtheorem{background}[theorem]{Background input}
\newtheorem{openproblem}[theorem]{Remaining task}

\newcommand{\Z}{\mathbb Z}
\newcommand{\1}{\mathbf 1}
\newcommand{\Hm}{H_m}
\newcommand{\Bm}{B_m}
\newcommand{\Hmx}{H_m^\times}
\newcommand{\Um}{\mathcal U_m}
\newcommand{\Em}{\widehat E_t}
\newcommand{\Gt}{\widehat G_t}
\newcommand{\Om}{\Omega_m}
\newcommand{\Dt}{D_m^\circ}

\title{Current theorem-order manuscript for the resonant \(d=5\) program:\\
local closure, visibility theorems, and the remaining frontier}
\author{}
\date{March 27, 2026}

\begin{document}
\maketitle

\begin{abstract}
This manuscript consolidates the current \(d=5\) proof program into a single paper-style \LaTeX{} document.
The purpose is \emph{not} to claim a finished proof for every resonant modulus.
Instead, the goal is to state in one place exactly what is already at theorem order, what now follows as a clean corollary, and what remains open.

The closed part of the program is the local symbolic machine for the promoted-plus one-line family
\[
P_m^+(t)=P_m^+\cup\{A_{t,3}\}, \qquad r=(m-1)/2,
\]
together with the corrected actual-entry reduction for the visibility problem
\[
\Hm \to \Bm.
\]
In the safe resonant regime \(m=3M\ge 27\) with \(5\nmid m\), the current theorem chain proves global visibility for the classes
\[
t=r-1,\qquad t=r,\qquad t=1,
\]
and also for the generic-late class when \(M\equiv 2 \pmod 5\).
For the exactness problem, the already-proved double-top law further shows that \(\Omega_m\)-leak is available for
\[
t=r-1,\qquad t=1,\qquad \text{generic late},
\]
but \emph{not} for \(t=r\).
Hence the proved exactness-ready families are currently \(t=r-1\), \(t=1\), and generic late with \(M\equiv 2 \pmod 5\), all conditional only on the remaining base-splice problem.

The remaining frontier is now explicit:
the unresolved generic-late arithmetic classes are \(M\equiv 1,3,4 \pmod 5\), the resonant base permutation \(B_m\) still needs a theorem-order single-cycle existence argument, and the branch \(5\mid m\) is not yet covered by the present global visibility theorems.
\end{abstract}

\tableofcontents

\section{Purpose and status convention}

The present document is a consolidation of the current \(d=5\) proof chain.
It is written so that the reader does not have to reconstruct the argument from dozens of short notes.
At the same time, it is important not to overstate the status.

We therefore distinguish three levels.

\begin{enumerate}[label=\textup{(\roman*)},leftmargin=2.8em]
\item \textbf{Theorem order.} A statement is already proved in the March 26--27 chain, with an explicit symbolic or arithmetic proof packet.
\item \textbf{Derived corollary.} A statement is not yet written as a separate note, but it follows immediately by combining theorem-order ingredients that are already closed.
\item \textbf{Frontier.} The statement is strongly audited or heuristically supported, but not yet promoted to a proof.
\end{enumerate}

The present manuscript records only theorem-order statements and derived corollaries as formal theorems.
The frontier is isolated in the final section.

\section{Global setting}

\subsection{The ambient graph-theoretic problem}

\begin{definition}[directed torus]
For an odd modulus \(m>2\), let
\[
D_5(m)=\operatorname{Cay}\big((\Z/m\Z)^5,\{e_0,e_1,e_2,e_3,e_4\}\big).
\]
A Hamilton decomposition of \(D_5(m)\) is a partition of its arc set into five arc-disjoint directed Hamilton cycles.
\end{definition}

The older March 24--25 package already closes the nonresonant packet \(3\nmid m\), the small packets \(m=5,7,9\), and the front-end surgery reductions.
The unresolved core is the resonant pure color-1 problem in the range \(m=3M\).
The present paper is about that resonant core.

\subsection{The promoted-plus one-line family}

Fix \(m\) odd, \(3\mid m\), and write
\[
r=\frac{m-1}{2}.
\]
The control family is
\[
P_m^+
=
\{
A_{1,4},\, C_{r,4},\, A_{1,3},\, C_{3,3},\, C_{r-1,3},\, A_{r,3},\, C_{r+1,3},\, A_{r+1,3}
\}.
\]
For each \(t\in\{0,\dots,m-1\}\) define the one-line extension
\[
P_m^+(t)=P_m^+\cup\{A_{t,3}\}.
\]

The color-1 section return on the \(\Sigma=0\) section will be denoted by
\[
R_t=h_1^m\big|_{\Sigma=0},
\]
and \(R_*=R_{P_m^+}\) denotes the control return.

\subsection{Section coordinates}

A \(\Sigma=0\) section state will be written as
\[
z=(a,c,d,e),
\]
with the suppressed coordinate
\[
b \equiv 1-a-c-d-e \pmod m.
\]
The coordinates \(a,b,c,d,e\) are always read modulo \(m\).

The visibility problem is studied on the set
\[
\Hm=\{(a,0,0,e): a,e\in \Z/m\Z\},
\qquad
\Bm=\{(a,0,0,0): a\in \Z/m\Z\},
\]
and we write
\[
\Hmx=\Hm\setminus \Bm.
\]

The compact coordinate on \(\Hmx\) is
\[
A=a-\1[e=m-1].
\]

\subsection{Auxiliary local sets}

The local proof chain uses the following subsets of the section.

\[
\Um=\{(a,c,m-2,e): 0\le c\le m-3\}
\]
(the hinge strip),
\[
K_m=\{(a,m-2,m-2,e)\}
\]
(the corner),
\[
S_m^-=\{(a,c,m-1,e): e\le m-2,\ c\ne m-2\}
\]
(the top-side strip),
\[
C_m^{\mathrm{full}}=\{(a,m-2,m-1,e)\}
\]
(the full gate),
\[
\Dt=\{(a,c,m-1,m-1): c\ne m-2\}
\]
(the double-top interior),
and
\[
\Om
=
\{(a,c,m-1,m-1):
c\not\equiv 1 \pmod 3,\ 
a\not\equiv 2c \pmod 3\}.
\]

\section{Baseline inputs already closed before the present pass}

The current paper does not reprove the entire March 24--25 front-end and nonresonant package from scratch.
Instead, it treats the following items as baseline inputs, and reorganizes the resonant core built on top of them.

\begin{background}[already closed baseline package]
The March 24--25 bundle supplies the following theorem-order inputs.
\begin{enumerate}[label=\textup{(\arabic*)},leftmargin=2.8em]
\item The nonresonant packet \(3\nmid m\) is closed.
\item The small odd packets \(m=5,7,9\) and the required surgery reductions are closed.
\item The width-1 section-return theorem and the promoted-collar control packet are available.
\item The residual front-end reductions from the main odd-\(m\), \(d=5\) manuscript are available.
\end{enumerate}
\end{background}

\begin{remark}
The purpose of the present manuscript is to absorb the \emph{new} resonant proof chain into one theorem ledger and to make the remaining open tasks explicit.
The March 24--25 baseline files are included in the archive accompanying this paper.
\end{remark}

\section{The local compact machine}

This section records the parts of the resonant local theory that are now at theorem order.

\subsection{Support localization and the lower strip}

\begin{theorem}[support localization and lower-strip invisibility]
For every \(t\in\{0,\dots,m-1\}\), the added row \(A_{t,3}\) is invisible on the whole strip \(d\le m-3\).
Consequently, on that strip the one-block kinematics are independent of \(t\) and satisfy
\[
c' = c+1 \pmod m,
\qquad
d' = d+\1[c=m-2].
\]
In particular, if \(c=0\) and \(d\le m-3\), then the next return time to \(c=0\) is exactly \(m\) blocks.
\end{theorem}

\begin{proof}[Proof idea]
The support-localization prefix lemma shows that the extra line \(A_{t,3}\) can only be seen after the orbit has reached the top/hinge active locus.
Before that moment the one-line family and the control family coincide.
The explicit width-1 lower-strip law therefore transfers from the control packet to all \(P_m^+(t)\), and the return time computation is then immediate from the deterministic \(c\)-advance.
\end{proof}

\subsection{The anomaly-free compact \(c=0\) return}

\begin{theorem}[compact \(c=0\) return law]
Let
\[
A=a-\1[e=m-1], \qquad
\lambda_d(A)=3+\1[A=1-d]-\1[A=0]-2\1[A=2-d]-\1[A=3-d].
\]
Then for every \(t\), every \(0\le d\le m-3\), and every state with \(c=0\), the first return after \(m\) blocks is
\[
R_t^m(a,0,d,e)=(A+\kappa^+,0,d+1,e^+),
\]
where
\[
e^+=e+\lambda_d(A), \qquad
\kappa^+=\1[e^+=m-1].
\]
Equivalently, the earlier low-band cocycle and the old top-band ``eight anomalies'' collapse into one anomaly-free compact law in the coordinate \(A\).
\end{theorem}

\begin{proof}[Proof idea]
The lower-strip theorem gives a full \(m\)-block excursion with no visibility of \(A_{t,3}\).
The old promoted-control low-band cocycle can therefore be upgraded to all one-line families.
The compact coordinate \(A=a-\1[e=m-1]\) is exactly the coordinate in which the apparent top-band anomalies disappear, because \(A\) is conserved through the \(c=0\) cycle and the final correction is absorbed into \(\kappa^+\).
\end{proof}

\subsection{Hinge conveyor, corner burst, and actual first top hits}

\begin{theorem}[hinge conveyor]
On the hinge strip \(\Um\), the induced one-block map is \(t\)-independent away from four boundary fibres and has the conveyor form
\[
(a,c,m-2,e)\mapsto (a+1,c+1,m-2,e).
\]
Consequently every orbit in \(\Um\) reaches the corner \(K_m\) in finitely many blocks.
\end{theorem}

\begin{theorem}[corner state reached from \(\Hmx\)]
Let
\[
z_0=(a_0,0,0,e_0)\in \Hmx,\qquad
A=a_0-\1[e_0=m-1].
\]
Immediately before the top episode begins, the orbit reaches the corner state
\[
z_C=(a_C,m-2,m-2,e_C),
\]
with
\[
e_C=e_0+\sigma(A),\qquad
a_C=A-2+\kappa_C,\qquad
\kappa_C=\1[e_C=m-1],
\]
where
\[
\sigma(A)
=
m-6+\1[A=0]-\1[A=2]+\1[A=3]+2\1[A=4]-\1[A=m-1].
\]
Equivalently,
\[
\sigma(A)=
\begin{cases}
m-4,& A=4,\\
m-5,& A\in\{0,3\},\\
m-6,& A\notin\{0,2,3,4,m-1\},\\
m-7,& A\in\{2,m-1\}.
\end{cases}
\]
\end{theorem}

\begin{theorem}[explicit actual-entry atlas]
Let
\[
\Phi_t:\Hmx\to \Em
\]
be the first-top-hit map.
For
\[
A=a_0-\1[e_0=m-1],\qquad
\epsilon=e_0+\sigma(A),\qquad
\kappa=\1[\epsilon=m-1],
\]
the actual first top hit \(\Phi_t(z_0)\) is given by Table~\ref{tab:actual-entry}.
Moreover:
\begin{enumerate}[label=\textup{(\arabic*)},leftmargin=2.8em]
\item \(\Em\) is supported on exactly two adjacent columns;
\item the side column is generated by the single compact residue \(A=4\);
\item the main-column row sizes are \(m-3,4,m-3\) at \(e=m-5,m-6,m-7\) and \(m-1\) elsewhere;
\item \(\Phi_t:\Hmx\to \Em\) is bijective.
\end{enumerate}
\end{theorem}

\begin{table}[t]
\centering
\caption{The actual-entry map \(\Phi_t\). The entries record \((a,c,e)\) at the first top hit.}
\label{tab:actual-entry}
\begin{tabular}{@{}p{0.20\textwidth}p{0.34\textwidth}p{0.34\textwidth}@{}}
\toprule
Class of \(t\) & Main branch \(A\ne 4\) & Anomaly branch \(A=4\) \\
\midrule
\(\{0,1,2\}\) & \((A,2,\epsilon)\) & \((5,1,\epsilon)\) \\
\(\{3\}\) & \((A+1,1,\epsilon)\) & \((4,2,\epsilon)\) \\
\(\{r-1\}\) & \((A-1,3,\epsilon)\) & \((4,2,\epsilon)\) \\
\(\{r\}\) & \((A+\kappa,1,\epsilon)\) & \((5+\kappa,0,\epsilon)\) \\
\(\{r+1\}\) & \((A+1,1,\epsilon)\) & \((6,0,\epsilon)\) \\
generic late & \((A,3,\epsilon)\) & \((5,2,\epsilon)\) \\
\bottomrule
\end{tabular}
\end{table}

\subsection{Top-row laws and the double-top packet}

\begin{theorem}[top-side transducer for \(t\ne 0\)]
On the strip \(S_m^-\), the one-block map is a finite class-dependent transducer.
Only the sigma-layers
\[
\{0,1,2,3,r-1,r,r+1,t\}
\]
can change the color-1 choice.
For \(t\ne 0\), the resulting transducer is explicit and produces the five arithmetic classes already isolated in the earlier scans:
\(t=1\), \(t=2\), \(t=r-1\), \(t\in\{3,r,r+1\}\), and generic late.
\end{theorem}

\begin{theorem}[the \(t=0\) global compact law]
Let \(E\) denote the explicit top-row early-cascade map proved for \(t=0\).
Then
\[
R_0 = E \quad\text{on }\{d=m-1\},
\qquad
R_0 = R_* \quad\text{on }\{d<m-1\}.
\]
Thus \(t=0\) is not a new global machine; it is an explicit top reset followed by the control continuation.
\end{theorem}

\begin{theorem}[double-top and \(\Omega_m\) law]
On the double-top interior \(\Dt\), the control family has an explicit choice word, and the whole one-line family \(P_m^+(t)\) admits an exact class-by-class symbolic law.
Restricted to the pocket \(\Om\), the behaviour is:
\begin{enumerate}[label=\textup{(\arabic*)},leftmargin=2.8em]
\item \(t=0,2\): immediate exit from \(\Om\);
\item \(t=1\): quarter-retention, hence eventual leak;
\item \(t=r-1\): half-retention, hence eventual leak;
\item \(t\in\{3,r,r+1\}\): \(\Om\) is invariant;
\item generic late: quarter-retention, hence eventual leak.
\end{enumerate}
\end{theorem}

\begin{remark}
Theorem 4.7 is the key reason that \(t=r\) is visibility-closed but not yet an exactness-ready family: on \(\Om\) the branch \(t=r\) belongs to the invariant class.
By contrast, \(t=1\), \(t=r-1\), and generic late are all \(\Om\)-leaking classes.
\end{remark}

\subsection{Corrected visibility reduction}

\begin{theorem}[corrected actual-entry reduction]
Let \(F_t:\Hm\to \Hm\) be the next-return map induced by \(R_t\).
Then the visibility problem
\[
\Hm \to \Bm
\]
is equivalent to the acyclicity of the finite induced map
\[
\Gt:\Em\to \Em\cup \Bm.
\]
More precisely, \(F_t\) has a cycle disjoint from \(\Bm\) if and only if \(\Gt\) has a cycle disjoint from \(\Bm\).
\end{theorem}

\begin{proof}[Proof idea]
Every non-base excursion from \(\Hm\) must pass through the corner and therefore through the actual first top-hit set \(\Em\).
Theorem 4.5 identifies those entries explicitly and bijectively.
After the top episode, the already-proved universal bridge brings the orbit either back to \(\Em\) or into \(\Bm\).
Thus the full visibility problem is exactly the cycle problem for the compact induced map \(\Gt\).
\end{proof}

\section{Global visibility theorems}

From this point on, the local mechanics are closed.
The remaining work is arithmetic and lives on the reduced compact map \(\Gt\).

\subsection{The already-closed classes \(t=r-1\) and \(t=r\)}

\begin{theorem}[visibility for \(t=r-1\)]
Assume
\[
m=3M\ge 27,\qquad 5\nmid m,\qquad t=r-1.
\]
Then the corrected actual-entry map \(\Gt\) is acyclic.
Equivalently, every orbit of \(F_t\) in \(\Hm\) meets \(\Bm\).
\end{theorem}

\begin{proof}[Proof idea]
The compact source law for \(t=r-1\) reduces the problem to a single \(A=0\) sweep.
The ordinary source map has one sweep survivor set \(\{2,7,12,17\}\), which moves after one sweep to \(\{5,10,15,20\}\), and none of those values survives a second sweep.
Hence the reduced source map has no non-base cycles, and therefore neither does \(\Gt\).
\end{proof}

\begin{theorem}[visibility for \(t=r\)]
Assume
\[
m=3M\ge 27,\qquad 5\nmid m,\qquad t=r.
\]
Then the corrected actual-entry map \(\Gt\) is acyclic.
Equivalently, every orbit of \(F_t\) in \(\Hm\) meets \(\Bm\).
\end{theorem}

\begin{proof}[Proof idea]
The compact source law for \(t=r\) again reduces the problem to a one-dimensional arithmetic sweep.
The active-source complement of the special window is exactly \(\{2,7,12\}\), and a full sweep shifts \(e\) by \(+3\).
Those survivors therefore move to \(\{5,10,15\}\), and for \(m\ge 33\) they do not meet the interface values.
The remaining case \(m=27\) is closed by an explicit short chain.
Thus the reduced source map is acyclic.
\end{proof}

\subsection{Active-source laws for \(t=1\) and generic late}

The next two results are the bridge from the local compact machine to the final arithmetic closure.

\begin{theorem}[explicit compact source law for \(t=1\) and generic late]
Assume
\[
m=3M\ge 27,\qquad 5\nmid m,
\]
and let \(t=1\) or let \(t\) be generic late
\[
t\notin \{0,1,2,3,r-1,r,r+1\}.
\]
Then the compact next-return on
\[
\Hmx=\{(A,e): A\in \Z/m\Z,\ e\in\{1,\dots,m-1\}\}
\]
has the following form.
\begin{enumerate}[label=\textup{(\arabic*)},leftmargin=2.8em]
\item On the ordinary branch \(e\ne s(A)\),
\[
e' = e+\rho(A),\qquad \rho(A)=m-5+\1[A=0]+2\1[A=4].
\]
\item If \(A\not\equiv 1 \pmod 3\), then \(A'=A\).
\item On the active class \(A\equiv 1 \pmod 3\), the ordinary branch is an explicit permutation \(p_t\).
\item On the special branch \(e=s(A)\), every state lands in \(\Bm\), except for the three interface outputs
\[
(m-1,1),\qquad (2,1),\qquad (3,m-1).
\]
\end{enumerate}
\end{theorem}

\begin{remark}
This theorem is the point at which the old ``source law gap'' disappears.
After it, the visibility problem is a cycle problem for an explicit one-dimensional map.
\end{remark}

\begin{theorem}[explicit active permutation theorem]
Let
\[
j=j(A)\in \Z_M
\]
be the active coordinate defined by
\[
A+2+6j\equiv 0 \pmod m.
\]
Then on the active ordinary branch, both for \(t=1\) and for generic late, the active source permutation has the form
\[
j' = j-r(j)-1 \pmod M,
\]
where \(r(j)\) is an explicit piecewise-constant gate remainder.
For \(t=1\) the remainder is controlled by \(M\bmod 4\); for generic late it is controlled by \(M\bmod 5\).
Thus the active permutation is an explicit finite interval-translation map with a short anomaly window.
\end{theorem}

\subsection{The \(t=1\) cycle theorem}

\begin{theorem}[the \(t=1\) active permutation cycle theorem]
Assume
\[
m=3M\ge 27,\qquad 5\nmid m,\qquad t=1.
\]
Then the active permutation has the following cycle structure.
\begin{enumerate}[label=\textup{(\arabic*)},leftmargin=2.8em]
\item If \(M\not\equiv 1 \pmod 3\), then it is a single \(M\)-cycle.
\item If \(M\equiv 1 \pmod 3\), write \(M=3q+1\). Then the cycle lengths are
\[
q,\qquad q,\qquad q+1.
\]
The unique cycle containing \(A=4\) is the one of length \(q+1\).
\end{enumerate}
\end{theorem}

\begin{proof}[Proof idea]
The explicit \(j\)-map becomes a finite column automaton.
When \(M=4q+1\), it is a height-4 column system plus one anomaly point; when \(M=4q+3\), it becomes a height-2 column system plus one anomaly point.
In both cases the base motion is a translation on \(\Z_q\) or \(\Z_{2q+1}\), and the cycle count reduces to \(\gcd(3,\cdot)\).
The anomaly contributes exactly one extra point to the unique base cycle containing \(A=4\).
\end{proof}

\begin{corollary}[visibility for \(t=1\)]
Assume
\[
m=3M\ge 27,\qquad 5\nmid m,\qquad t=1.
\]
Then the corrected actual-entry map \(\Gt\) is acyclic.
Equivalently,
\[
\Hm \to \Bm
\]
visibility holds for \(t=1\).
\end{corollary}

\begin{proof}
The compact source law is Theorem 5.3.
The active cycle decomposition is Theorem 5.5.
The conditional arithmetic closure from the earlier visibility sweep note then becomes unconditional:
fixed classes drift by \(-4\) or \(-5\) and therefore hit their special rows because \(5\nmid m\); interface states lie on those fixed classes and also terminate in base; and on each active cycle the one-sweep shift
\[
S_C=-5|C|+2\1[4\in C]
\]
has \(\gcd(S_C,m)\in\{1,3\}\), while the forbidden set meets every sweep orbit.
Hence no non-base source cycle exists, and therefore \(\Gt\) is acyclic.
\end{proof}

\subsection{The generic-late class}

\begin{theorem}[generic late, \(M\equiv 2\pmod 5\)]
Assume
\[
m=3M\ge 27,\qquad 5\nmid m,
\]
and let \(t\) be generic late. Suppose
\[
M\equiv 2 \pmod 5,\qquad M=5q+2.
\]
Then the active permutation has the following cycle structure.
\begin{enumerate}[label=\textup{(\arabic*)},leftmargin=2.8em]
\item If \(3\nmid q\), then it is a single \(M\)-cycle.
\item If \(3\mid q\), then the cycle lengths are
\[
Q,\qquad Q+1,\qquad Q+1,
\qquad
Q=\frac{M-2}{3}.
\]
\end{enumerate}
\end{theorem}

\begin{proof}[Proof idea]
In the congruence class \(M=5q+2\), the generic active map becomes a height-5 column automaton plus two anomaly points.
On the base set one sees the translation \(b\mapsto b-3\) on \(\Z_q\), with exactly two edges replaced by two-step paths through the anomaly states.
The cycle count is then again controlled by \(\gcd(3,q)\), and the two anomaly points are inserted on the residue-0 and residue-2 cycles.
\end{proof}

\begin{corollary}[generic-late visibility when \(M\equiv 2 \pmod 5\)]
Assume
\[
m=3M\ge 27,\qquad 5\nmid m,
\]
let \(t\) be generic late, and suppose \(M\equiv 2 \pmod 5\).
Then the corrected actual-entry map \(\Gt\) is acyclic.
Equivalently,
\[
\Hm \to \Bm
\]
visibility holds for generic late in this congruence class.
\end{corollary}

\begin{proof}
Combine the explicit compact source law (Theorem 5.3), the explicit active map (Theorem 5.4), the cycle theorem above, and the same one-sweep arithmetic closure used in the proof of Corollary 5.6.
The fixed classes and interface states terminate because \(5\nmid m\), while each active-cycle sweep has no non-base cycles.
\end{proof}

\section{What the current theorem chain implies for exactness}

The exactness problem is stronger than visibility because it also requires the double-top pocket to leak and the base permutation to splice.

\begin{proposition}[local-global exactness criterion]
Assume for a given class \(t\) that:
\begin{enumerate}[label=\textup{(\arabic*)},leftmargin=2.8em]
\item the pocket \(\Om\) has no closed \(R_t\)-cycle (\(\Omega_m\)-leak);
\item every non-pocket orbit meets \(\Hm\);
\item every \(\Hm\)-cycle meets \(\Bm\) (visibility);
\item the next-return permutation on \(\Bm\) is a single cycle (base splice).
\end{enumerate}
Then the full color-1 section return is exact, namely it consists of one cycle of length \(m^4\).
\end{proposition}

\begin{proof}
Any section cycle either lies inside \(\Om\) or meets \(\Hm\).
The first case is ruled out by \(\Omega\)-leak.
The second case meets \(\Bm\) by visibility, and once the cycle meets \(\Bm\), the single-cycle assumption on the base next-return forces all such intersections to belong to one and the same global section cycle.
Hence there is only one section cycle.
\end{proof}

\begin{corollary}[exactness-ready classes in the current theorem package]
In the safe resonant regime \(m=3M\ge 27\), \(5\nmid m\), the present theorem chain reduces the pure color-1 exactness problem to the base-splice problem for the following classes:
\[
t=r-1,\qquad t=1,\qquad \text{generic late with } M\equiv 2 \pmod 5.
\]
By contrast, the class \(t=r\) is visibility-closed but \(\Omega_m\)-invariant, so it is not exactness-ready.
\end{corollary}

\begin{proof}
Visibility is proved in Theorems 5.1, 5.2, Corollary 5.6, and Corollary 5.8.
The double-top theorem shows \(\Omega\)-leak for \(t=r-1\), \(t=1\), and generic late, but not for \(t=r\).
Apply Proposition 6.1.
\end{proof}

\begin{remark}
The exact computational witnesses found earlier -- such as \((m,t)=(27,1)\), \((33,4)\), \((39,1)\), and the staged \(m=51\) cases -- fit exactly this picture: once leak and visibility are both in place, the only unresolved theorem-order burden is the single-cycle base splice.
\end{remark}

\section{What is still missing from a full resonant \(d=5\) proof}

The present paper organizes the remaining work into a very short list.

\begin{openproblem}[generic late, remaining mod-5 classes]
For generic late, the arithmetic visibility theorem is still open in the congruence classes
\[
M\equiv 1,\ 3,\ 4 \pmod 5.
\]
The active source permutation is explicit in all three cases; what remains is the cycle theorem for the corresponding interval-translation maps.
\end{openproblem}

\begin{openproblem}[base single-cycle existence]
In the exactness-ready classes
\[
t=r-1,\qquad t=1,\qquad \text{generic late with }M\equiv 2 \pmod 5,
\]
the remaining pure color-1 burden is to prove that the induced base permutation on \(\Bm\) has a single cycle, or to prove an equivalent splice certificate.
This is the current splice problem.
\end{openproblem}

\begin{openproblem}[the branch \(5\mid m\)]
All global visibility theorems in the present manuscript are proved in the safe regime \(5\nmid m\).
Some top-row classes are known to have trapped cycles when \(5\mid m\), while \(t=1\) and \(t=r\) remain plausible candidates.
A dedicated arithmetic source analysis is still needed there.
\end{openproblem}

\begin{openproblem}[final all-color assembly and \(d=\)prime]
Once the remaining resonant pure color-1 splice problem is solved, the staged color \(0/2\) repairs and the broader \(d=\)prime architecture should be rewritten in the same compact-machine language.
The present manuscript already indicates the likely template:
support localization, compact return, actual-entry atlas, visibility, and splice.
\end{openproblem}

\section{A one-page theorem ledger}

Table~\ref{tab:ledger} is the shortest honest summary of the current state.

\begin{longtable}{@{}p{0.39\textwidth}p{0.22\textwidth}p{0.31\textwidth}@{}}
\caption{Current theorem ledger for the resonant \(d=5\) core.}
\label{tab:ledger}\\
\toprule
Item & Status & Role \\
\midrule
\endfirsthead
\toprule
Item & Status & Role \\
\midrule
\endhead
Nonresonant packet \(3\nmid m\), small odd packet \(m=5,7,9\) & baseline input & already closed before the present pass \\
Support localization and lower-strip kinematics & theorem order & removes \(t\)-dependence below the top two rows \\
Compact \(c=0\) return law in the coordinate \(A=a-\1[e=m-1]\) & theorem order & unifies low-band and top-band return laws \\
Hinge conveyor and corner state formula & theorem order & funnels every non-base excursion to one corner state \\
Actual-entry atlas \(\Phi_t:\Hmx\to \Em\) & theorem order & turns visibility into a finite actual-entry problem \\
Top-row transducer for \(t\ne 0\) & theorem order & closes the local top mechanics \\
\(t=0\) global compact law & theorem order & shows \(t=0\) is explicit reset plus control continuation \\
Double-top / \(\Omega_m\) law & theorem order & isolates leak versus invariant-wall classes \\
Corrected visibility reduction \(\Gt:\Em\to \Em\cup\Bm\) & theorem order & finite model for \(\Hm\to\Bm\) \\
Visibility for \(t=r-1\) & theorem order & arithmetic closure in a splice-safe class \\
Visibility for \(t=r\) & theorem order & arithmetic closure in a non-leaking class \\
Compact source law for \(t=1\) and generic late & theorem order & reduces global visibility to a 1D source map \\
\(t=1\) active permutation cycle theorem & theorem order & closes the \(t=1\) active cycle count \\
Generic late, \(M\equiv 2\pmod 5\) cycle theorem & theorem order & closes one full generic mod-5 branch \\
Visibility for \(t=1\) & derived corollary & now unconditional by combining theorem-order ingredients \\
Visibility for generic late, \(M\equiv 2\pmod 5\) & derived corollary & same \\
Generic late, \(M\equiv 1,3,4\pmod 5\) & open & remaining visibility arithmetic frontier \\
Base single-cycle existence on \(\Bm\) & open & remaining exactness obstruction in the closed visibility classes \\
\(5\mid m\) resonant branch & open & outside the present safe-regime theorems \\
\bottomrule
\end{longtable}

\section{Bottom line}

The problem is no longer ``understand the whole \(m^4\)-state dynamics''.
It has been reduced to the compact scheme
\[
\text{compact return}
\longrightarrow
\text{actual-entry atlas}
\longrightarrow
\text{finite visibility map}
\longrightarrow
\text{base splice}.
\]

What is already proved is substantial:

\begin{itemize}[leftmargin=2.4em]
\item the local compact machine is closed;
\item visibility is proved for \(t=r-1\), \(t=r\), \(t=1\), and generic late with \(M\equiv 2\pmod 5\);
\item exactness is therefore reduced to the base-splice problem in the three exactness-ready families \(t=r-1\), \(t=1\), and generic late with \(M\equiv 2\pmod 5\).
\end{itemize}

What remains is equally explicit:

\begin{itemize}[leftmargin=2.4em]
\item generic late mod-5 classes \(1,3,4\);
\item the resonant base single-cycle theorem;
\item the safe-to-\(5\mid m\) extension and the final all-color assembly.
\end{itemize}

That is the present theorem-order status of the \(d=5\) program.

\end{document}
